\documentclass[UTF8,a4paper,12pt]{ctexart}

\usepackage{geometry}
\usepackage{fontspec}
\usepackage{xeCJK}
\usepackage{amsmath}

\geometry{
  left=2.5cm,
  right=2.5cm,
  top=2.5cm,
  bottom=2.5cm
}

% ------------------------------------------------------------
% Windows 中文字体测试
% 如果你已经导入 Windows fonts.zip，这里应该直接命中真实字体。
% 如果你配置了 fontconfig alias，也可以测试 alias 是否生效。
% ------------------------------------------------------------
\setCJKmainfont{SimSun}[
  AutoFakeBold=true,
  AutoFakeSlant=true
]

\setCJKfamilyfont{song}{SimSun}[
  AutoFakeBold=true,
  AutoFakeSlant=true
]

\setCJKfamilyfont{fangsong}{FangSong}[
  AutoFakeBold=true,
  AutoFakeSlant=true
]

\setCJKfamilyfont{hei}{SimHei}[
  AutoFakeBold=true,
  AutoFakeSlant=true
]

\setCJKfamilyfont{kai}{KaiTi}[
  AutoFakeBold=true,
  AutoFakeSlant=true
]

\newcommand{\songtiwin}{\CJKfamily{song}}
\newcommand{\fangsongwin}{\CJKfamily{fangsong}}
\newcommand{\heitiwin}{\CJKfamily{hei}}
\newcommand{\kaitiwin}{\CJKfamily{kai}}

\title{中文字体测试文档}
\author{TeXDock / ShareLaTeX 字体环境测试}
\date{\today}

\begin{document}

\maketitle

\section{基础中文测试}

这是默认正文。当前正文主字体应为宋体 SimSun。

中文测试：你好，世界。  
标点测试：，。！？；：“”‘’（）《》——……  
繁体测试：臺灣、資料庫、軟體工程、雲端服務。  
英文混排测试：This is a XeLaTeX Chinese font test.

\section{宋体测试}

{\songtiwin
这是宋体 SimSun 测试段落。

宋体通常用于正文。这里测试普通中文、英文 ABC、数字 1234567890，以及数学公式旁边的文字。

加粗测试：\textbf{这是宋体加粗测试。}

倾斜测试：\textit{这是宋体倾斜测试。}
}

\section{仿宋测试}

{\fangsongwin
这是仿宋 FangSong 测试段落。

仿宋常用于公文、标题说明或正式文档中的特定段落。

加粗测试：\textbf{这是仿宋加粗测试。}

倾斜测试：\textit{这是仿宋倾斜测试。}
}

\section{黑体测试}

{\heitiwin
这是黑体 SimHei 测试段落。

黑体通常用于标题、强调文字或小节名称。

加粗测试：\textbf{这是黑体加粗测试。}
}

\section{楷体测试}

{\kaitiwin
这是楷体 KaiTi 测试段落。

楷体通常用于引文、说明、题词或示例文本。

加粗测试：\textbf{这是楷体加粗测试。}
}

\section{中英文与数学混排测试}

这是一个中英文混排段落：The quick brown fox jumps over the lazy dog。

数学公式测试：

\[
E = mc^2
\]

\[
\sum_{i=1}^{n} i = \frac{n(n+1)}{2}
\]

中文公式说明：上面的公式用于测试数学环境是否正常，中文字体是否和数学字体共存。

\section{结论}

如果本文档能够正常编译，并且 PDF 中可以看到宋体、仿宋、黑体、楷体的不同字形，说明中文字体、fontconfig alias、XeLaTeX 和 ctex 基本工作正常。

\end{document}
